\begin{abstract}
    We introduce the \emph{central distribution}, a probability distribution that represents a set or hierarchy of
    probability distributions by minimizing the worst-case Kullback--Leibler (KL)
    divergence to the admissible models.
    This provides a principled method for converting partial knowledge and set-valued uncertainty into a single
    probability measure suitable for Bayesian inference.

    For finite sets of distributions, we establish existence, uniqueness, and minimax characterization results, and
    show that central distributions arise as solutions to convex--concave saddle-point problems.
    For parametric models, we derive a sufficient fixed-point condition for priors whose induced marginals
    are central, revealing a formal connection with Bernardo's reference prior methodology.

    We extend the framework to hierarchical uncertainty structures and propose an asymptotic notion of centrality for
    infinite sequences.
    We also derive approximation bounds that enable practical construction of approximate central distributions,
    allowing Bayesian inference without explicit knowledge of the exact solution.

    Finally, we demonstrate how uncertainty about relevant feature subsets can be represented hierarchically and
    aggregated into a single predictive model.
    These results suggest a general information-theoretic framework for representing partial knowledge, constructing
    noninformative priors, and reasoning under model uncertainty.
\end{abstract}
